\subsubsection{Robustness of the generated-regressor correction across DGP specifications}

To test whether the Robinson correction improves the plug-in estimator under alternative data generating processes, we vary two features of the DGP: the distribution of $x$ and the dependence of $\nu$ on $\eta$. Throughout, $x \perp \eta$ unconditionally, $\gamma = \beta = 2$, and treatment is $s = \mathbf{1}\{q > 0\}$ with $q = \gamma x + \eta$. The nonparametric component is approximated by a cubic B-spline with $k_n = \lfloor n^{1/3}\rceil$ knots on fixed support $[-2.8,\,2.8]$. We run $R = 500$ replications at $n = 50{,}000$ for each combination, comparing the Robinson estimator (which includes a shifted-basis correction for the generated regressor) to the naive estimator (single-step OLS with no correction). The key statistic is the variance ratio $\mathrm{Var}(\hat\alpha_{\mathrm{Rob}}) / \mathrm{Var}(\hat\alpha_{\mathrm{Naive}})$: values below 1 indicate the correction helps.

\paragraph{Specifications tested.}
For $x$: (i) $N(0,1)$; (ii) $t(3)/\sqrt{3}$ (heavy tails, unit variance); (iii) $\mathrm{Exp}(1) - 1$ (right-skewed, unit variance); (iv) bimodal mixture $\frac{1}{2}N(-1.5,\,0.25) + \frac{1}{2}N(1.5,\,0.25)$, rescaled to unit variance.
For $\nu$: (i) baseline $\nu \mid a \sim N(3e^a/(1+e^a),\,1)$; (ii) $\nu \mid \eta \sim N(\eta^2,\,1)$ (quadratic, $h' = 2\eta$); (iii) $\nu \mid \eta \sim N(e^{\eta/2},\,1)$ (exponential, $h'$ varies by factor $\sim\!16$ across support); (iv) $\nu \sim N(0,1)$ independent of $\eta$ ($h_{\mathrm{Con}} = 0$); (v) $\nu \mid \eta \sim N(2\eta,\,1)$ (linear, $h' = \mathrm{const}$).

\begin{table}[h]
\centering
\caption{Robinson vs.\ Naive across DGP specifications ($n = 50{,}000$, $\gamma = 2$). The exponential-$x$ rows are unreliable (only 50--113 of 500 replications survived the outlier filter) and are included for completeness.}
\label{tab:dgp_sweep}
\small
\begin{tabular}{ll r r r r r r}
\toprule
$x$ dist & $\nu$ spec & $n_{\mathrm{ok}}$ & $\mathrm{Cor}_{\mathrm{Tr}}$ & Bias$_R$ & Bias$_N$ & Std$_R$ / Std$_N$ & VarRatio \\
\midrule
normal  & baseline     & 500 & $-0.45$ & 0.015 & 0.015 & 0.022\,/\,0.022 & 1.01 \\
normal  & $\eta^2$     & 500 & $-0.45$ & 0.399 & 0.408 & 0.077\,/\,0.075 & 1.06 \\
normal  & $e^{\eta/2}$ & 500 & $-0.45$ & 0.307 & 0.320 & 0.052\,/\,0.052 & 1.03 \\
normal  & independent  & 500 & $-0.45$ & $-0.001$ & $-0.001$ & 0.018\,/\,0.018 & 1.00 \\
normal  & $2\eta$      & 500 & $-0.45$ & 0.146 & 0.146 & 0.024\,/\,0.024 & 1.01 \\
\addlinespace
$t(3)$  & baseline     & 500 & $-0.37$ & 0.006 & 0.006 & 0.022\,/\,0.022 & 1.01 \\
$t(3)$  & $\eta^2$     & 500 & $-0.36$ & 0.297 & 0.310 & 0.089\,/\,0.086 & 1.07 \\
$t(3)$  & $e^{\eta/2}$ & 500 & $-0.36$ & 0.226 & 0.241 & 0.056\,/\,0.053 & 1.10 \\
$t(3)$  & independent  & 500 & $-0.37$ & 0.000 & 0.000 & 0.021\,/\,0.020 & 1.01 \\
$t(3)$  & $2\eta$      & 500 & $-0.37$ & 0.097 & 0.097 & 0.027\,/\,0.027 & 1.01 \\
\addlinespace
bimodal & baseline     & 500 & $-0.49$ & 0.017 & 0.018 & 0.030\,/\,0.030 & 1.04 \\
bimodal & $\eta^2$     & 500 & $-0.49$ & 0.618 & 0.631 & 0.100\,/\,0.095 & 1.11 \\
bimodal & $e^{\eta/2}$ & 500 & $-0.49$ & 0.375 & 0.374 & 0.076\,/\,0.074 & 1.07 \\
bimodal & independent  & 500 & $-0.49$ & $-0.001$ & 0.000 & 0.027\,/\,0.027 & 1.02 \\
bimodal & $2\eta$      & 500 & $-0.49$ & 0.175 & 0.176 & 0.034\,/\,0.034 & 1.01 \\
\addlinespace
exp     & baseline     & 101 & $-0.46$ & 0.116 & 0.071 & 0.525\,/\,0.529 & 0.98 \\
exp     & $\eta^2$     &  50 & $-0.46$ & 0.519 & 0.073 & 0.429\,/\,0.474 & 0.82 \\
exp     & $e^{\eta/2}$ &  64 & $-0.47$ & 0.272 & 0.256 & 0.513\,/\,0.518 & 0.98 \\
exp     & independent  & 113 & $-0.46$ & 0.001 & $-0.037$ & 0.503\,/\,0.506 & 0.99 \\
exp     & $2\eta$      &  97 & $-0.46$ & 0.259 & 0.182 & 0.521\,/\,0.525 & 0.99 \\
\bottomrule
\end{tabular}
\end{table}

Several patterns emerge from Table~\ref{tab:dgp_sweep}. First, no combination produces a variance ratio below 1.0 among the reliable specifications (normal, $t(3)$, bimodal). The correction never reduces variance. Second, the correction is most harmful when $h'$ varies strongly: the $\eta^2$ and $e^{\eta/2}$ specifications produce the largest variance ratios (up to 1.11 for bimodal $\times$ $\eta^2$). This is the opposite of what one might expect---more $h'$ variation should, in principle, make the correction term less collinear with $x$ and thus more informative. Instead, the additional regressor introduces multicollinearity noise that exceeds the bias it removes. Third, the linear specification ($\nu \mid \eta \sim N(2\eta,1)$, constant $h'$) serves as a sanity check: the correction term is exactly proportional to $x$, so VarRatio $\approx 1.00$ everywhere, confirming it adds no information. Fourth, the bimodal distribution produces the strongest selection-induced correlation ($\mathrm{Cor} = -0.49$) and also the largest variance penalties from the correction, while $t(3)$ has the weakest ($\mathrm{Cor} = -0.36$) and correspondingly smaller penalties.

The exponential-$x$ results are unreliable: the right-skewed distribution causes the plug-in estimator to fail in the majority of replications ($\mathrm{Std} \sim 0.5$, only 10--23\% of runs survive the outlier filter). The apparent VarRatio $< 1$ in the exp/$\eta^2$ cell ($0.82$) is based on only 50 surviving replications and is not robust.

In summary, across 20 DGP combinations spanning symmetric and skewed $x$ distributions, constant and highly variable $h'$, and selection-induced correlations ranging from $-0.36$ to $-0.49$, the Robinson correction provides no variance reduction for the plug-in estimator. It is at best irrelevant and consistently slightly harmful, with the penalty increasing when $h'$ has more variation or when the within-subsample correlation between $x$ and $\hat\eta$ is stronger.
